
We observed significant heterogeneity in the estimated effects of school closures on SARS-CoV-2 infections and 
COVID-19 deaths. In this section, we explore potential sources of this heterogeneity by examining the 
relationships between the estimated effects of school closures and key country-specific factors, 
including demographics and the non-pharmaceutical interventions in place during the study period. 

First, we evaluated the correlation between the estimated effects of school closures on infections and deaths, and the following factors:
\begin{itemize}
    \item proportion of individuals aged under 15 years old, 
    \item proportion of individuals aged 70 years old and above,
    \item proportion of enrolled students in the population,
    \item total duration of school closure,
    \item average stringency of PHSMs other than school closures implemented during school closures.
\end{itemize}

Table \ref{tab:correlation} presents the correlation between each of these factors and the 
estimated effect of school closures on infections and deaths. 

\begin{table}[!ht]
    \begin{center}
    \begin{tabular}{|p{9cm}|p{3.5cm}|p{3.5cm}|}
    \hline
    \textbf{Correlation} & \textbf{Intervention effect on SARS-CoV-2 infections} & \textbf{Intervention effect on COVID-19 deaths} \\
    \hline
    N weeks of schools closures & 0.43 & 0.33 \\
    Average Oxford Stringency index during school closures & -0.18 & -0.20 \\
    \% of population under 15 years old (year 2020) & 0.46 & 0.31 \\
    \% of population enrolled in schools (year 2020) & 0.44 & 0.34 \\
    \% of population above 70 years old (year 2020) & -0.54 & -0.39 \\
    \hline
    \end{tabular}
\end{center}
    % \end{adjustbox}
    \caption{Correlation between various factors and the estimated intervention effect on SARS-CoV-2 infections and COVID-19 deaths}
    \label{tab:correlation}
\end{table}

Based on these univariate correlation coefficients, we observe that school closures may have more beneficial effects on 
infections and deaths in settings that exhibit one or more of the following characteristics:
\begin{enumerate}
    \item a large proportion of individuals aged under 15, 
    \item a large proportion of enrolled students,
    \item less stringent PHSMs (other than school closures) during school periods
    \item a small proportion of individuals aged 70 and above
    \item a longer total duration of school closures
\end{enumerate}

It is important to note that calculating the correlation coefficients constitute a crude approach that was only used as a first step in our explorations. The following sections present how we calculated the variables associated with each of these factors and 
their detailed relationships with the estimated effect of school closures on infections and deaths, respectively.

\subsection{Proportion of individuals aged under 15 years old}
The proportion of individuals aged under 15 years old was calculated from the population age distributions
reported by the United Nations' Population Division for year 2020. 

Figure \ref{fig:hetero_kids} presents how the estimated effects of school closures on infections and deaths varied 
with the proportion of individuals aged under 15.

\begin{figure}[!ht]
    \begin{center}
    \includegraphics[width=0.7\textwidth]{../../../../user/rragonnet/remote_run_outputs/heterogeneity_plots/hetero_prop_kids.pdf}
    \end{center}
    \caption{Relationship between school closure impact on infections (top panel) and deaths (bottom panel), and 
    proportion of individuals aged under 15 years old. Intervention effects are presented as median (dots) and ranges (bars), 
    with colours indicating the continents to which the countries belong.} 
    \label{fig:hetero_kids}
\end{figure}


\subsection{Proportion of individuals aged 70 years old and above}
The proportion of individuals aged 70 years old and above was calculated from the population age distributions
reported by the United Nations' Population Division for year 2020. 

Figure \ref{fig:hetero_elderly} presents how the estimated effects of school closures on infections and deaths varied 
with the proportion of individuals aged 70 and above.

\begin{figure}[!ht]
    \begin{center}
    \includegraphics[width=0.7\textwidth]{../../../../user/rragonnet/remote_run_outputs/heterogeneity_plots/hetero_prop_elderly.pdf}
    \end{center}
    \caption{Relationship between school closure impact on infections (top panel) and deaths (bottom panel), and 
    proportion of individuals aged 70 years old and over. Intervention effects are presented as median (dots) and ranges (bars), 
    with colours indicating the continents to which the countries belong.} 
    \label{fig:hetero_elderly}
\end{figure}


\subsection{Proportion of enrolled students}
For each country, the total number of enrolled students was obtained from the UNESCO database on school closures
for year 2020. The proportion of enrolled students was then obtained by dividing the number of enrolled students 
by the total population size reported by the United Nations' Population Division for year 2020.

Figure \ref{fig:hetero_students} presents how the estimated effects of school closures on infections and deaths varied 
with the proportion of enrolled students in the population.

\begin{figure}[!ht]
    \begin{center}
    \includegraphics[width=0.7\textwidth]{../../../../user/rragonnet/remote_run_outputs/heterogeneity_plots/hetero_prop_students.pdf}
    \end{center}
    \caption{Relationship between school closure impact on infections (top panel) and deaths (bottom panel), and 
    proportion of students in the population. Intervention effects are presented as median (dots) and ranges (bars), 
    with colours indicating the continents to which the countries belong.} 
    \label{fig:hetero_students}
\end{figure}


\subsection{Total duration of school closures}
The total duration of school closures was calculated using the UNESCO school closure database. 
The total number of school closure weeks was calculated by adding the number of weeks fully closed, 
and the number of weeks partially closed multiplied by 0.3 (our midpoint estimate for assumed attendance 
fraction during partial closures).

Figure \ref{fig:hetero_duration} presents how the estimated effects of school closures on infections and deaths varied 
with the total duration of school closures. Note that the bottom panel of Figure \ref{fig:hetero_duration} 
presents the same data as Figure 2 in the main text.

\begin{figure}[!ht]
    \begin{center}
    \includegraphics[width=0.7\textwidth]{../../../../user/rragonnet/remote_run_outputs/heterogeneity_plots/hetero_n_weeks_closed.pdf}
    \end{center}
    \caption{Relationship between school closure impact on infections (top panel) and deaths (bottom panel), and 
    total duration of school closures. Intervention effects are presented as median (dots) and ranges (bars), 
    with colours indicating the continents to which the countries belong.} 
    \label{fig:hetero_duration}
\end{figure}


\subsection{Average PHSM stringency index}
We used stringency index data from The Oxford Covid-19 Government Response Tracker (OxCGRT) to inform our calculations
of the average stringency index. Namely, we used the same approach as that used by 
OxCGRT to compute their \textit{StringencyIndex\_Average} index, but we excluded the school closure
component from the calculations. Thus, the resulting average stringency index reflects the total 
extent to which PHSMs other than school closures were implemented in each country. A higher 
value of the index indicates a higher PHSM stringency level.

Figure \ref{fig:hetero_stringency} presents how the estimated effects of school closures on infections and deaths varied 
with the average stringency index for PHSMs other than school closures.

\begin{figure}[!ht]
    \begin{center}
    \includegraphics[width=0.7\textwidth]{../../../../user/rragonnet/remote_run_outputs/heterogeneity_plots/hetero_stringency.pdf}
    \end{center}
    \caption{Relationship between school closure impact on infections (top panel) and deaths (bottom panel), and 
    average stringency of PHSMs other than school closures. Intervention effects are presented as median (dots) and ranges (bars), 
    with colours indicating the continents to which the countries belong.} 
    \label{fig:hetero_stringency}
\end{figure}


\subsection{Multivariate analysis}
In order to account for potential confounding, we performed a multivariate analysis to further explore
the heterogeneity observed in the estimated effects of school closures on infections and deaths.

First, we noted that some of the factors investigated in the previous sections were dependent. In particular,
three factors were directly relevant to the age-structure of the population and therefore highly dependent: proportion under 15 years old, proportion 70 years and above, and proportion of students. We 
used a causal directed acyclic graph (cDAG) to better represent these dependencies in the context of our intervention
effects estimations (Figure \ref{fig:cDAG}).

\begin{figure}[!ht]
    \begin{center}
    \includegraphics[width=0.5\textwidth]{../../../../user/rragonnet/remote_run_outputs/heterogeneity_plots/cDAG.jpg}
    \end{center}
    \caption{Simple causal directed acyclic graph (cDAG) for the effect of school closures on SARS-Cov-2 
    infections and distal outcome COVID-19 deaths. This cDAG informed the construction of a simple linear model 
    examining this effect. Green circle=main exposure of interest, Blue circles=outcomes of interest, Boxes=adjusted confounding covariates. PHSMs = public health and social measures.} 
    \label{fig:cDAG}
\end{figure}

The cDAG supports including three factors in a multilinear model: the population age structure, 
the total duration of school closures, and the stringency of PHSMs other than 
school closures. Among the factors related to population age structure, the correlation coefficients suggest that 
the proportion of individuals aged 70 and above may have the strongest relationship with the estimated effects of 
school closures. Consequently, we incorporated the following three variables into our multilinear regression analysis:
\begin{itemize}
    \item proportion of individuals aged 70 years old and above,
    \item total duration of school closure,
    \item average stringency of PHSMs other than school closures implemented during school closures.
\end{itemize}

Figure \ref{fig:multireg} presents the results of the multilinear regression.

\begin{figure}[!ht]
    \centering
    \subfigure[]{
        \includegraphics[width=0.8\textwidth]{../../../../user/rragonnet/remote_run_outputs/heterogeneity_plots/linreg_infections.png}
        \label{fig:reg_inf}
    }
    \\ % Line break to stack the images vertically
    \subfigure[]{
        \includegraphics[width=0.8\textwidth]{../../../../user/rragonnet/remote_run_outputs/heterogeneity_plots/linreg_deaths.png}
        \label{fig:reg_death}
    }
    \caption{Results of the multilinear regression exploring the relationship between different population factors and the effects of school closures
    on infections (panel a) and deaths (panel b). The factors are: the proportion of individuals aged 70 years old and above (\textit{prop\_elderly}), 
    the average stringency of PHSMs other than school closures implemented during the school closures (\textit{stringency}), 
    and the total duration of school closures (\textit{n\_weeks\_closed}).
    }
    \label{fig:multireg}
\end{figure}

The proportion of elderly was found to be the only statistically significant predictor (p-value $\leq 0.05$) of the effect of school closures
on infections and deaths. This suggests that settings exhibiting a higher proportion of elderly individuals are less likely to benefit from
school closures in terms of both number of infections and deaths. The other two factors were not found to be significantly associated with the outcome variables.